\chapter{Teenagers and adolescent s mental health}
\index{Teenagers and adolescent s mental health}

\medskip
\noindent\textit{Educational reference; always seek professional advice for individual care.}

\section*{Definition and Overview}
Adolescence is a period of enormous physical, emotional, and social change, and mental health during these years matters for the teenage years themselves and for the whole of adult life. Adolescent mental health refers to the emotional, psychological, and social wellbeing of young people aged roughly ten to nineteen: how they feel, think, relate to others, cope with stress, and manage the challenges of growing up. Most teenagers experience ups and downs, and periods of moodiness, worry, and withdrawal are normal as they form their own identity and face exams, friendships, and first relationships. A mental-health problem exists when these difficulties are severe, persistent, or out of proportion, or when they interfere with school, family life, friendships, or basic daily functioning. The common problems of this age include anxiety, depression, eating disorders, attention deficit hyperactivity disorder (ADHD), behavioural difficulties, and, less commonly, psychosis. Mental illness in adolescence is common, often begins in these years, and is very treatable when recognised early and approached without judgement. This chapter explains what adolescent mental health looks like, which signs should prompt help-seeking, how problems are assessed and treated, and where young people and families can find support.

\section*{Epidemiology}
Mental-health difficulties are among the most common health problems of the teenage years. Around one in seven young people worldwide experiences a mental disorder, and most adult mental illness begins before the age of twenty-five. Anxiety is the most common problem, followed by depression, with rates rising sharply through the teenage years, particularly in girls. Large local surveys find that a substantial minority of young people experience a mental-health disorder in any given year, and that many do not receive professional help, which is why awareness and early action matter so much. Some groups face higher rates, including Aboriginal and Torres Strait Islander young people, those in out-of-home care, LGBTIQ+ young people, and those who have experienced trauma, bullying, or poverty. Eating disorders most often begin in adolescence, and suicide is one of the leading causes of death among young people, which underlines why talk of self-harm or suicidal thoughts must always be taken seriously. The encouraging side is that the teenage brain is highly responsive to treatment: the earlier a problem is recognised, the better the outcome, and most young people recover with the right help.

\section*{Aetiology and Pathophysiology}
Teenage mental-health problems arise from the interaction of brain development, genetics, and the world the young person lives in. No single cause exists; rather, several influences combine.
\begin{itemize}
    \item \textbf{Brain development:} the emotional centres of the brain mature earlier than the prefrontal cortex that controls impulses and planning, so teenagers are primed to feel intensely and act on emotion before judgement catches up
    \item \textbf{Hormones and puberty:} the hormonal surges of adolescence reshape mood, sleep, and appetite, and the timing of puberty can affect wellbeing
    \item \textbf{Genetics and family history:} anxiety, depression, and other conditions run in families, and a family history of mental illness raises vulnerability
    \item \textbf{Social and academic pressure:} exams, school transitions, and performance expectations can drive stress and perfectionism
    \item \textbf{Peer relationships and bullying:} friendship problems, cyberbullying, and social rejection are among the strongest risk factors for depression and anxiety in this age group
    \item \textbf{Family conflict and adversity:} parental conflict, divorce, abuse, neglect, grief, and financial stress shape a young person's sense of safety and self-worth
    \item \textbf{Trauma:} violence, abuse, or loss can trigger depression, anxiety, and post-traumatic stress
    \item \textbf{Sleep disruption:} staying up late, screens in the bedroom, and early school starts produce chronic sleep deprivation, which is strongly linked to low mood and anxiety
    \item \textbf{Substance use:} alcohol and other drugs can trigger, mimic, or worsen mental-health problems, and are sometimes used to self-medicate distress
    \item \textbf{Social media:} heavy use is associated with poorer sleep, body-image concerns, and comparison-driven distress, though it is also a source of connection for many
\end{itemize}

\section*{Clinical Features}
The signs of a mental-health problem in a teenager are changes from their usual self, and the changes are often first noticed by parents, teachers, or friends.

\begin{itemize}
    \item \textbf{Persistent low mood or irritability:} sadness, grumpiness, or a short temper lasting for weeks, rather than a few bad days
    \item \textbf{Loss of interest:} dropping out of sport, hobbies, and activities that used to bring pleasure, and seeming to no longer care
    \item \textbf{Worry that is out of proportion:} constant worry about school, health, or social acceptance, avoidance of situations, panic attacks, or physical symptoms such as a racing heart
    \item \textbf{Withdrawal:} isolating in the bedroom, declining invitations, and pulling back from friends and family
    \item \textbf{Sleep and appetite changes:} difficulty falling asleep, waking at night, oversleeping, or changes in eating and weight
    \item \textbf{Declining school performance:} falling grades, missed classes, difficulty concentrating, and lost motivation
    \item \textbf{Risk-taking:} reckless behaviour, drinking, drug use, or dangerous activity that is out of character
    \item \textbf{Physical symptoms:} persistent headaches, stomach aches, or tiredness with no clear physical cause
    \item \textbf{Self-harm and suicidal thoughts:} cutting, overdosing, or other self-injury, or talking about hopelessness, death, or suicide, which must always be taken seriously
\end{itemize}

Not every teenager shows all of these, and a single sign on its own may be a phase, but several signs together, or any sign lasting more than a couple of weeks, is worth a check-in.

\section*{Differential Diagnosis}
A professional assessment separates genuine mental-health conditions from normal development, physical illness, and other causes of changed behaviour.

\begin{itemize}
    \item \textbf{Normal adolescent moodiness:} short-lived grumpiness and withdrawal that respond to support are part of growing up; the key differences are persistence, severity, and impact on functioning
    \item \textbf{Depression:} persistent low mood, loss of interest, sleep and appetite change, and poor concentration lasting at least two weeks
    \item \textbf{Anxiety disorders:} generalised anxiety, panic, social anxiety, and separation anxiety have worry and avoidance at their core and are the most common mental-health conditions of adolescence
    \item \textbf{ADHD:} long-standing inattention, impulsivity, and restlessness that began in childhood and interfere with school and daily life
    \item \textbf{Autism:} social communication differences, sensory sensitivities, and a need for routine that can be mistaken for withdrawal or odd behaviour
    \item \textbf{Eating disorders:} restrictive eating, binge eating, purging, intense fear of weight gain, and an overvalued sense of shape and weight
    \item \textbf{Substance use:} intoxication, withdrawal, and drug-related mood swings can look like mental illness, and the two often coexist
    \item \textbf{Physical illness:} thyroid problems, glandular fever, and iron deficiency cause tiredness, low mood, and poor concentration and should be excluded
    \item \textbf{Psychosis:} hearing voices, unusual fixed beliefs, paranoia, and a decline in function are red flags requiring urgent specialist assessment
\end{itemize}

\section*{When to Seek Medical Care}
See a GP if changes in mood, behaviour, sleep, or school function last for more than a couple of weeks, if they are getting worse, or if they are affecting daily life. A GP can check for physical causes, prepare a mental health care plan that gives access to subsidised psychology sessions, and refer to a psychologist or psychiatrist. If there is any concern about self-harm or suicidal thoughts, seek help immediately and do not wait: asking directly about suicide does not put the idea into someone's head. For non-urgent support, call a a local youth crisis service, a a local mental-health support service, or a a a local crisis service. Aboriginal and Torres Strait Islander young people can call a local culturally appropriate crisis service for a culturally safe yarn with a crisis supporter.

\textbf{Contact your local emergency services for:}
\begin{itemize}
    \item Any suicide attempt or self-harm needing medical attention
    \item A plan or immediate intent to end their life, or to harm someone else
    \item Behaviour that puts the young person in immediate danger
    \item Severe confusion, agitation, or hallucinations
\end{itemize}

\section*{Diagnosis and Investigations}
The assessment of a teenager's mental health is a careful process that relies on conversation, observation, and information from school and family, rather than on a single test.

\begin{itemize}
    \item \textbf{Detailed history:} the clinician asks the young person and their family about the onset and pattern of symptoms, sleep, appetite, school, friendships, and any stressful events
    \item \textbf{Mental-state examination:} an interview that assesses mood, thinking, memory, concentration, and whether there are thoughts of harm to self or others
    \item \textbf{Questionnaires:} validated tools such as the PHQ-A for depression and GAD-7 for anxiety help measure severity and track progress, but are only aids, not diagnoses
    \item \textbf{Information from school:} with consent, teachers can describe how the young person is functioning in class and among peers over time
    \item \textbf{Physical review:} a medical check and, if indicated, blood tests to exclude thyroid problems, anaemia, and glandular fever
    \item \textbf{Substance-use screening:} a confidential, non-judgemental conversation about alcohol and drug use
    \item \textbf{Specialist referral:} complex, severe, or risky presentations are referred to a child and adolescent psychiatrist or a public adolescent mental-health service
\end{itemize}

\section*{Management}
Most teenage mental-health problems respond well to treatment, and the treatment is tailored to the condition, its severity, and the young person's own goals.

\begin{itemize}
    \item \textbf{Psychoeducation:} helping the young person and family understand the condition, its triggers, and what recovery looks like reduces fear and stigma and builds cooperation
    \item \textbf{Cognitive behavioural therapy (CBT):} the first-line therapy for depression and anxiety; it teaches young people to recognise unhelpful thoughts and behaviours and to build coping skills
    \item \textbf{Interpersonal therapy:} an evidence-based therapy for adolescent depression that focuses on relationships and conflict
    \item \textbf{Family therapy:} involving the family improves communication and support, and is particularly valuable in younger adolescents and eating disorders
    \item \textbf{Medication:} for moderate-to-severe depression and anxiety, selective serotonin reuptake inhibitors (SSRIs) may be added under specialist supervision, with careful monitoring in the first weeks
    \item \textbf{School support:} working with the school on adjustments, tutoring, counselling, and exam arrangements keeps the young person connected to learning
    \item \textbf{Early-psychosis services:} dedicated teams assess and treat psychosis-like symptoms quickly, because early treatment greatly improves the long-term outlook
    \item \textbf{Lifestyle support:} regular sleep, daily movement, and reduced screen time are active treatment for mood and anxiety
    \item \textbf{Online and phone support:} programs such as ReachOut, eheadspace, and a local mental-health support service offer accessible, low-stigma options for young people
\end{itemize}

\section*{Complications}
\begin{itemize}
    \item \textbf{Social isolation:} withdrawal and stigma erode friendships and family connection at the age when belonging matters most
    \item \textbf{Substance use:} young people often turn to alcohol and drugs to numb distress, which worsens the underlying problem
    \item \textbf{Self-harm:} cutting and other self-injury are common ways of coping with unbearable feelings, and carry real physical risks
    \item \textbf{Suicide:} the most serious consequence, and one of the leading causes of death in young people
    \item \textbf{Chronicity:} without treatment, adolescent depression and anxiety often persist or recur into adulthood
    \item \textbf{Family strain:} a young person's mental illness affects the whole household, and parents and siblings need support too
\end{itemize}

\section*{Prognosis and Outcomes}
The prognosis for adolescent mental-health problems is genuinely good. Most young people with depression and anxiety recover fully with appropriate treatment, and the teenage brain's flexibility means that skills learned in therapy take hold well. Early intervention makes the biggest difference: the shorter the time between symptoms starting and treatment beginning, the better the outcome, which is why the affected region invests in services such as headspace. Some conditions, such as ADHD and some mood and psychotic disorders, are more persistent and are managed over years rather than cured, but even these respond well and allow young people to complete school, study, work, and build relationships. Relapses can occur, particularly in depression, and part of treatment is learning to recognise early warning signs and seek help early the next time. With a combination of evidence-based therapy, family support, school involvement, and, where needed, medication, the large majority of teenagers recover and move into adult life with resilience and hope.

\section*{Prevention and Self-Care}
\begin{itemize}
    \item \textbf{Keep conversations about feelings normal:} ask your teenager directly how they are, then listen without fixing
    \item \textbf{Protect sleep:} a consistent bedtime, no screens in the bedroom at night, and about eight to ten hours of sleep make a measurable difference to mood
    \item \textbf{Encourage daily movement and time outdoors,} which are as effective for mild low mood as many other interventions
    \item \textbf{Limit and make time for social media,} and talk openly about comparison, body image, and online bullying
    \item \textbf{Keep family routines such as shared meals,} which build connection and give an early window into how a teenager is really going
    \item \textbf{Take any talk of self-harm or suicide seriously,} ask about it directly, and seek help the same day
    \item \textbf{Encourage a mental health care plan from a GP,} which gives access to subsidised psychology sessions
    \item \textbf{Look after your own wellbeing as a parent,} and remember that asking for help for your teenager is a strength, not a failure
\end{itemize}

\section*{Sources and Further Reading}
The following organisations provide reliable information and support for young people and their families.

\begin{itemize}
    \item headspace. \textit{Mental health support for young Australians.} https://headspace.org.au/
    \item Beyond Blue. \textit{Young people's mental health and wellbeing.} https://www.beyondblue.org.au/
    \item Kids Helpline. \textit{Free, private counselling for young people.} https://kidshelpline.com.au/
    \item Lifeline. \textit{Crisis support, 24 hours a day.} https://www.lifeline.org.au/
    \item healthdirect Australia. \textit{Teenage mental health: signs and getting help.} https://www.healthdirect.gov.au/
    \item Raising Children Network. \textit{Teenage mental health for parents.} https://raisingchildren.net.au/
    \item Orygen. \textit{Youth mental health research and resources.} https://www.orygen.org.au/
\end{itemize}

\clearpage
